\documentclass[12pt]{article}
\usepackage{fancyhdr}
\pagestyle{fancy}
\fancyhf{}
\fancyfoot[C]{\thepage}
\fancyfoot[L]{\scriptsize CC BY-NC-SA 4.0 \textbullet\ Leo Liang}
\renewcommand{\headrulewidth}{0pt}
\usepackage[margin=1in]{geometry}
\usepackage{amsmath, amssymb}
\usepackage{array, booktabs}
\usepackage{pgffor}
\parindent=0pt
\parskip=1em
\begin{document}
\pagestyle{fancy}
\begin{center}
  {\Large\bfseries Test 1 Topics}\\[4pt]
  Speaking Math \;|\; Sets \;|\; Functions \\[6pt]
  1.5-hour test
\end{center}
\hrule\bigskip

\subsection*{General Notes}

The homework and your question compasses are a good study guide. If you understand all the homework problem, you will most likely do well. 

You do \textbf{not} need to remember the exact definitions of anything. We will \textbf{not} ask you a question like ``What is the definition of \_\_\_\_\_\_\_\_". Just understand the intuition, be able to explain it in words, and be able to deeply logically reason about it. If needed, use the internet and AI (such as ChatGPT, Claude, Gemini, etc.) for help. You can ask something like ``Give me 3 sample problems for this topic \textbf{X}," ``Explain \textbf{X} to me more, I don't quite understand it,", ``Explain \textbf{X} in more simple terms. Use an analogy," ``How does \textbf{X} relate to [an idea from a different unit]," ``How did you solve this problem and why did you do this step \textbf{X}?" If you use AI, make sure to ask it \textbf{specific questions}. If you use AI,\textbf{ use it as a teacher, not your answer key}.

Manage your time wisely. Communicate with Leo concerns or problems you're facing :)

Good luck and have fun thinking.

\section*{Talking (Speaking Math)}
\begin{itemize}
  \item Expression vs.\ equation vs.\ formula
  \item Terms, coefficients, variables, constants, exponents, like terms
  \item Reading algebra aloud (``the quantity,'' powers, fractions, etc.)
  \item Writing expressions from words
\end{itemize}

\section*{Sets}
\begin{itemize}
  \item Sets \& membership ($\in$, $\notin$)
  \item Roster \& set-builder notation
  \item Number sets ($\mathbb{N}$, $\mathbb{Z}$, $\mathbb{R}$): which one are larger than the others? and why?
  \item Subsets ($\subseteq$): proper and improper subset
  \item Union, intersection, difference, complement
  \item Cardinality ($|A|$)
  \item Cartesian products ($A \times B$)
  \item Inclusion--exclusion
  \item Real-world applications
\end{itemize}

\section*{Functions}
\begin{itemize}
  \item The four representations (arrow, ordered pairs, table, graph)
  \item Function notation $f(x)$
  \item Domain, codomain, range: \quad be able to explain the similarities and differences between all the three in words. be prepared to be asked specific questions about them.
  \item Injective, surjective, bijective: \quad be able to explain the similarities and differences between all the three in words. be prepared to be asked specific questions about them.
  \item Proofs (injective, surjective, and bijective proof)
  \item Real-world applications: \quad think of examples like the Minecraft bedroom and variations of that
  \item How do we compare the sizes of infinite sets?
\end{itemize}
\textbf{Note:} Composition, inverses, the verticle-line test, and multivariable functions will \textbf{not} be on the test.

\begin{center}\rule{0.5\linewidth}{0.4pt}\end{center}
{\large\bfseries Part 2 --- New Material (Limits)}

\section*{Domain \& Range}
\begin{itemize}
  \item \textbf{Domain} --- avoid the two dangers: no dividing by zero, no square roots of negative values.
  \item \textbf{Range} --- set $y = f(x)$ and solve backwards for $x$; which outputs are actually reachable? Understand the ``story" that the equation $x = g(y)$ is telling you.
  \item Reading the domain and range off a graph
\end{itemize}
\section*{Limits --- The Big Ideas}
\begin{itemize}
  \item What is a limit?
  \item One-sided limits (from the left, from the right) and two-sided limits.
  \item ``Does not exist (DNE)'' vs.\ ``undefined'': \quad be able to explain the difference in words, and why a limit can exist where the function is undefined (a hole)
  \item Reading limits from a graph: open vs.\ filled circles, holes, and jumps
  \item Finding limits from a table of values
  \item Finding limits with algebra: try substituting first; if you get $\tfrac{\text{something}}{0}$, factor and cancel
  \item Infinite limits and vertical asymptotes ($\pm\infty$)
\end{itemize}
\section*{Limits at Infinity (End Behavior)}
\begin{itemize}
  \item What ``$\displaystyle\lim_{x\to\infty} f(x)$'' is asking, and how it shows up as a horizontal asymptote \textit{in the long term}. (Notice we say ``in the long term" because the range may still include the value of the horizontal asymptote as shown in class).
  \item Finding $\lim\limits_{x\to\infty} f(x)$ with the shortcut and algebraic approach.
  \item Know the \textbf{limit laws}: sum, difference, constant multiple, product, quotient, power.
\end{itemize}

\end{document}